\section{Chemická rovnováha}

\ce{a A + b B <=> c C + d D}

$Q = \frac{[C]^c[D]^d}{[A]^a[B]^b}$ - reakční kvocient

$\Delta = \Delta G^0 + RT\ln Q\\
K = \frac{[C]_r^c[D]_r^d}{[A]_r^a[B]_r^b}\\
\Delta G^0 = -RT\ln K$

V rovnováze platí Q = K.

$K = e^{-\frac{\Delta G^0}{RT}}$

Pro reakce v plynné fázi můžeme místo koncentrace použít tlak.

\ce{N2(g) + 3 H2(g) <=> 2 NH3(g)}

$K = \frac{p^2_{NH3}}{p_{N2}\ .\ p_{H2}^3}$

U heterogenních systémů je aktivita kondenzovaných fází rovna jedné.

\ce{3 Fe2O3(s) + H2(g) <=> 2 Fe3O4(s) + H2O(g)}

$K = \frac{a^3_{Fe3O4\ .\ a_{H2O}}}{a_{Fe2O3}^3\ .\ a_{H2}} = \frac{p_{H2O}}{p_{H2}}$

\subsection{Le-Chatelierův princip}

\textit{Systém, který je v rovnováze, reaguje na každou změnu tak, aby tuto změnu potlačil.}

\begin{itemize}
	\item Exotermní vs endotermní reakce
	\item Reakce s plyny, ovlivnění tlakem (výroba amoniaku)
	\begin{itemize}
		\item \ce{H2(g) + I2(g) -> 2 HI(g)}
		\item Přídavek inertního plynu rovnováhu neovlivní -- parciální tlak
	\end{itemize}
	\item Vliv koncentrace -- esterifikace
\end{itemize}